\documentclass{thesisclass}
% Based on thesisclass.cls of Timo Rohrberg, 2009
% ----------------------------------------------------------------
% Thesis - Main document
% ----------------------------------------------------------------
% No empty in-between chapters, remove this line if desired for double-page printing
\let\cleardoublepage\clearpage

%% -------------------------------
%% |  Information for PDF file   |
%% -------------------------------
\hypersetup{
 pdfauthor={Not set},
 pdftitle={Not set},
 pdfsubject={Not set},
 pdfkeywords={Not set}
}


%% ---------------------------------
%% | Information about the thesis  |
%% ---------------------------------

\newcommand{\myname}{Fangyu Xu}
\newcommand{\mytitle}{Spectral and Decomposition Tracking for Rendering Heterogeneous
Volumes}
\newcommand{\myinstitute}{Institut für Visualisierung und Datenanalyse,\\ Lehrstuhl für Computergrafik}

\newcommand{\reviewerone}{?}
\newcommand{\reviewertwo}{?}
\newcommand{\advisor}{?}
\newcommand{\advisortwo}{?}

\newcommand{\timestart}{XX. Monat 20XX}
\newcommand{\timeend}{XX. Monat 20XX}
\newcommand{\submissiontime}{DD. MM. 20XX}


%% ---------------------------------
%% | ToDo Marker - only for draft! |
%% ---------------------------------
% Remove this section for final version!
\setlength{\marginparwidth}{20mm}

\newcommand{\margtodo}
{\marginpar{\textbf{\textcolor{red}{ToDo}}}{}}

\newcommand{\todo}[1]
{{\textbf{\textcolor{red}{(\margtodo{}#1)}}}{}}


%% --------------------------------
%% | Old Marker - only for draft! |
%% --------------------------------
% Remove this section for final version!
\newenvironment{deprecated}
{\begin{color}{gray}}
{\end{color}}


%% --------------------------------
%% | Settings for word separation |
%% --------------------------------
% Help for separation:
% In german package the following hints are additionally available:
% "- = Additional separation
% "| = Suppress ligation and possible separation (e.g. Schaf"|fell)
% "~ = Hyphenation without separation (e.g. bergauf und "~ab)
% "= = Hyphenation with separation before and after
% "" = Separation without a hyphenation (e.g. und/""oder)

% Describe separation hints here:
\hyphenation{
% Pro-to-koll-in-stan-zen
% Ma-na-ge-ment  Netz-werk-ele-men-ten
% Netz-werk Netz-werk-re-ser-vie-rung
% Netz-werk-adap-ter Fein-ju-stier-ung
% Da-ten-strom-spe-zi-fi-ka-tion Pa-ket-rumpf
% Kon-troll-in-stanz
}


%% ------------------------
%% |    Including files   |
%% ------------------------
% Only files listed here will be included!
% Userful command for partially translating the document (for bug-fixing e.g.)
\includeonly{%
titlepage,
content
}


%%%%%%%%%%%%%%%%%%%%%%%%%%%%%%%%%
%% Here, main documents begins %%
%%%%%%%%%%%%%%%%%%%%%%%%%%%%%%%%%
\begin{document}

% Remove the following line for German text
%\selectlanguage{ngerman}

\frontmatter
\pagenumbering{roman}
%% titlepage.tex
%%

% coordinates for the bg shape on the titlepage
\newcommand{\diameter}{20}
\newcommand{\xone}{-15}
\newcommand{\xtwo}{160}
\newcommand{\yone}{15}
\newcommand{\ytwo}{-253}

\begin{titlepage}
% bg shape
\begin{tikzpicture}[overlay]
\draw[color=gray]  
 		 (\xone mm, \yone mm)
  -- (\xtwo mm, \yone mm)
 arc (90:0:\diameter pt) 
  -- (\xtwo mm + \diameter pt , \ytwo mm) 
	-- (\xone mm + \diameter pt , \ytwo mm)
 arc (270:180:\diameter pt)
	-- (\xone mm, \yone mm);
\end{tikzpicture}
	\begin{textblock}{10}[0,0](4,2.5)
		\includegraphics[width=.3\textwidth]{logos/KITLogo_RGB.pdf}
	\end{textblock}
	\changefont{phv}{m}{n}	% helvetica	
	\vspace*{3.5cm}
	\begin{center}
		\Huge{\mytitle}
		\vspace*{2cm}\\
		\Large{
			Proseminar-Ausarbeitung von
		}\\
		\vspace*{1cm}
		\huge{\myname}\\
		\vspace*{1cm}
		\Large{
			An der Fakult\"at f\"ur Informatik
			\\
			\myinstitute
		}\\
		\vspace*{1cm}
		\Large{\today}
	\end{center}
	\vspace*{1cm}
%\Large{
%\begin{center}
%\begin{tabular}[ht]{l c l}
%  % Gutachter sind die Professoren, die die Arbeit bewerten. 
%  \iflanguage{english}{Reviewer}{Erstgutachter}: & \hfill  & \reviewerone\\
%  \iflanguage{english}{Second reviewer}{Zweitgutachter}: & \hfill  & \reviewertwo\\
%  \iflanguage{english}{Advisor}{Betreuender Mitarbeiter}: & \hfill  & \advisor\\
%  \iflanguage{english}{Second advisor}{Zweiter betreuender Mitarbeiter}: & \hfill  & \advisortwo\\
%  % Der zweite betreuende Mitarbeiter kann weggelassen werden. 
%\end{tabular}
%\end{center}
%}


\vspace{2cm}


\begin{textblock}{10}[0,0](4,16.8)
\tiny{ 
		{KIT -- Universität des Landes Baden-Württemberg und nationales Forschungszentrum der Helmholtz-Gesellschaft}
}
\end{textblock}

\begin{textblock}{10}[0,0](14,16.75)
\large{
	\textbf{www.kit.edu} 
}
\end{textblock}

\end{titlepage}

\blankpage


%% -------------------
%% |   Directories   |
%% -------------------
\tableofcontents
\blankpage


%% -----------------
%% |   Main part   |
%% -----------------
\mainmatter
\pagenumbering{arabic}
%% content.tex
%%

%% ==============================
\chapter{Introduction}
\label{ch:Introduction}
%% ==============================

\section{Introduction of Spectral and Decomposition Tracking}
In this paper, the authors focused on two more efficient unbiased rendering techniques, decomposition tracking and spectral tracking, for volume rendering in heterogeneous medium. Decomposition tracking is based on delta tracking, which requires doing a volume lookup at each step. In order to reduce such a large consumption of memory and computation, the authors used decomposition tracking to speed the generation of free-path up by splitting the medium into a homogeneous control and a heterogeneous residual. Then by adding a fictitious medium, the heterogeneous residual is homogenized. (...) The authors also contributed another technique, spectral tracking, which makes efficient tracking in chromatic media possible. (...)
\section{Existing Tracking Algorithms}
Introduce the existing tracking algorithms: Delta Tracking ...
Closed-form Tracing.
Delta Tracking. The main idea of delta tracking is to sample a heterogeneous volume by using the closed-form homogeneous-volume sampling. Delta tracking uses fictitious medium, which is represented as the null-collision coefficient $\mu_n\left(x\right)$, to create homogeneous medium. The fictitious medium has no impacts on the light transport, which means when a fictitious particle is hit, the light continues going in the same direction; when a real particle is hit, the light will be scattered. All these are done to allow (...).

\section{Radiative Transfer with Null Collisions}
Describe the RTE, delta tracking and weighted delta tracking.\\
differential form of radiative transfer equation: \[\left(\omega \cdot \bigtriangledown \right)L\left(x, \omega \right) = -\mu_t\left(x\right) L\left(x, \omega \right) + \mu_a\left(x\right) L_e\left(x, \omega \right) + \mu_s\left(x\right)  \int_{S^2} f_p\left(\omega, \bar{\omega} \right)\,d\bar{\omega}(1)\]


%% ==============
\chapter{Decomposition Tracking}
\label{ch:Content1}
%% ==============
Decomposition tracking presented by the authors is a volume-decomposition technique for accelerating free-path sampling. It decomposes the original medium into two components, one is homogeneous control component $\{ \mu_a^c, \mu_s^c\}$, the other is heterogeneous residual component $\{ \mu_a^r, \mu_s^r\}$.
Section 2.1(...) Section 2.2(...)

%% ===========================
\section{Analog Decomposition Tracking}
\label{ch:Content1:sec:Section1}
%% ===========================



%% ===========================
\section{Weighted Decomposition Tracking}
\label{ch:Content1:sec:Section2}
%% ===========================

\dots


%% content.tex
%%

%% ==============
\chapter{Spectral Tracking}
\label{ch:Content2}
%% ==============

\dots


%% ===========================
\section{Collision Probabilities}
\label{ch:Content2:sec:Section1}
%% ===========================

\dots


%% ===========================
\section{Spectral and Decomposition Tracking}
\label{ch:Content2:sec:Section2}
%% ===========================

\dots

%% ==============
\chapter{Conclusion}
\label{ch:Content1}
%% ==============



%% --------------------
%% |   Bibliography   |
%% --------------------
\cleardoublepage
\phantomsection
\addcontentsline{toc}{chapter}{\bibname}

\iflanguage{english}
{\bibliographystyle{IEEEtranSA}}	% english style
{\bibliographystyle{babalpha-fl}}	% german style
												  
% Use IEEEtran for numeric references
%\bibliographystyle{IEEEtranSA})

\bibliography{ausarbeitung}
\Erklaerung
\end{document}
